\documentclass[11pt]{article}
\usepackage{amsmath,amssymb,amsfonts,amsthm,graphicx,booktabs,hyperref}
\usepackage[margin=2.4cm]{geometry}
\usepackage{microtype}
\hypersetup{colorlinks=true,linkcolor=blue,citecolor=blue,urlcolor=blue}
\newtheorem{theorem}{Theorem}
\newtheorem{lemma}{Lemma}
\newtheorem{proposition}{Proposition}
\newtheorem{corollary}{Corollary}
\theoremstyle{definition}
\newtheorem{definition}{Definition}
\newtheorem{numresult}{Numerical result}
\title{Protective Delay is a Different Loop:\\ Two Institutional Channels}
\author{Your Name}
\date{\today}
\begin{document}
\maketitle

\begin{abstract}
The companion core destabilises when a scarcity-mobilising institution raises effort in response to a filtered deficit. That claim is specific to the sign of the delayed gain $C_Z$. A protective institution that tracks a decreasing quota map $E_{\mathrm{cap}}(Z)$ at the same Candidate~A equilibrium has $C_Z<0$, loop gain $0.080$, and no positive root of the Hopf cubic: it is delay-independently stable. Replacing $C_Z$ by $-C_Z$ at fixed modulus leaves $H$ unchanged and shifts $\tau$ by $\pi/\omega$; the resulting linearisation is not the quota law. The two-channel system with delays $(\tau_m,\tau_p)$ has characteristic identity $P(\lambda)=L(\lambda)\bigl(C_m e^{-\lambda\tau_m}+C_p e^{-\lambda\tau_p}\bigr)$. If the weighted loop gain is strictly less than one, there is no delay-induced Hopf for any $(\tau_m,\tau_p)$. Destabilisation by short delay is confined to the mobilising summand.
\end{abstract}

\section{Mobilising and protective linearisations}
\label{sec:claim}

The companion~\cite{paperI} proves, for the effort-bounded three-state core
\begin{equation}
\dot N=S(N)-qEN,\qquad
\dot Z=\tau_m^{-1}\bigl(\Phi(qEN-S(N))-Z\bigr),\qquad
\dot E=\Bigl(1-\frac{E}{E_{\max}}\Bigr)F(E,Z(t-\tau)),
\label{eq:core}
\end{equation}
that a scarcity-mobilising $F$ with
\[
F=\eta E\Bigl(\frac{Z_\tau}{\Delta_{\mathrm{ref}}}-\frac{E}{E_{\max}}\Bigr)+\delta_0\frac{Z_\tau}{Z_{\mathrm{ref}}+Z_\tau}
\]
produces a pair of local Hopfs at Candidate~A, interval-certified as
\[
\tau_-\in[3.6661490142739,\,3.6661490142743],\qquad
\tau_+\in[150.3584773101408,\,150.3584773101421].
\]
The linearisation about the interior equilibrium is
\[
\dot x=A_N x+A_E y,\quad
\dot z=B_N x+B_E y-d z,\quad
\dot y=C_E y+C_Z z(t-\tau),
\]
with $C_Z=+1.785$ at gated Candidate~A. The sign of $C_Z$ distinguishes the institutional channel. The opposite channel and the two-channel interpolation are treated below.

\section{Protective tracking}

\begin{definition}[Quota-tracking protection]
\label{def:prot}
A protective institution is the law
\begin{equation}
\dot E=\Bigl(1-\frac{E}{E_{\max}}\Bigr)\eta_p\bigl(E_{\mathrm{cap}}(Z(t-\tau_p))-E\bigr),
\label{eq:prot}
\end{equation}
where $E_{\mathrm{cap}}$ is $C^2$, positive, and strictly decreasing. The calibration
\begin{equation}
E_{\mathrm{cap}}(Z)=E_0\frac{Z_{\mathrm{ref}}}{Z_{\mathrm{ref}}+Z},\qquad
E_0=E^*_A\frac{Z_{\mathrm{ref}}+\delta}{Z_{\mathrm{ref}}}
\label{eq:cap}
\end{equation}
places the unique interior rest of~\eqref{eq:prot} at the companion's Candidate~A point $(N^*,Z^*,E^*)=(89.55188,\delta,2.08962)$, so the stock--memory block is identical and the only change is the effort law.
\end{definition}

At that rest $E_{\mathrm{cap}}(\delta)=E^*$ and
\begin{equation}
C_E=-\Bigl(1-\frac{E^*}{E_{\max}}\Bigr)\eta_p,\qquad
C_Z=\Bigl(1-\frac{E^*}{E_{\max}}\Bigr)\eta_p E_{\mathrm{cap}}'(\delta).
\label{eq:CZ}
\end{equation}
With $\eta_p=\eta_A=0.914$,
\[
C_E=-0.850336,\qquad C_Z=-1.661702.
\]
Both signs are those of a restoring quota, not of scarcity mobilisation.

\section{Delay-independent stability}

The companion's Hopf cubic $H(x)$ depends on $C_Z$ only through $C_Z^2$. A sign flip of $C_Z$ at fixed modulus does not change $H$. A genuine protective law changes the modulus as well.

\begin{theorem}[No delay-induced Hopf under quota tracking]
\label{thm:nohopf}
At the Candidate~A stock--memory linearisation and the gains~\eqref{eq:CZ}, the cubic $H$ has no positive root. Equivalently,
\[
\sup_{\omega>0}\bigl|L(i\omega,\tau_p)\bigr|
=\sup_{\omega>0}\left|\frac{C_Z\,L(i\omega)}{(i\omega-A_N)(i\omega+d)(i\omega-C_E)}\right|
=0.08011<1.
\]
The companion's loop-gain theorem then excludes every nonzero imaginary characteristic root for all $\tau_p\ge 0$. In addition, the undelayed characteristic system is Hurwitz and the zero-root condition is delay independent with $P(0)\ne0$ because $L(0)=0$. By continuous dependence of retarded characteristic roots, exponential stability therefore persists for every $\tau_p\ge0$.
\end{theorem}

\begin{proof}
The stock--memory coefficients $(A_N,A_E,B_N,B_E,d)$ are those of Candidate~A in~\cite{paperI}. Substituting~\eqref{eq:CZ} gives the monic cubic
\[
H(x)=x^3+c_2 x^2+c_1 x+c_0,
\]
with numerical coefficients
\[
c_2=A_N^2+d^2+C_E^2=0.76339,\qquad
c_1=0.028946,\qquad
c_0=9.278\times 10^{-6},
\]
all positive, and $c_2 c_1-c_0=0.02209>0$. Descartes' rule of signs therefore yields no positive real root, and the Routh--Hurwitz criterion confirms that all roots of $H$ have negative real parts when viewed as a polynomial in a complex variable $x$. Consequently $|P(i\omega)|=|C_Z|\,|L(i\omega)|$ has no solution $\omega>0$.

Equivalently, the loop gain $\Gamma(\omega)=|C_Z L(i\omega)/P(i\omega)|$ is continuous on $(0,\infty)$, tends to $0$ as $\omega\to 0$ and as $\omega\to\infty$, and, in the externally verified computation, attains maximum $0.08011$ near $\omega\approx 0.0583$, computed by setting $\Gamma'(\omega)=0$ on a logarithmically spaced grid and refining the critical point by Newton. In particular $\sup\Gamma<1$, so the small-gain criterion of~\cite{paperI} excludes every imaginary-axis root for all $\tau_p\ge 0$.
\end{proof}

\begin{numresult}[Protective spectrum at Candidate A]
\label{num:prot}
The characteristic quasi-polynomial of~\eqref{eq:core} with~\eqref{eq:prot} in place of the mobilising bracket has no imaginary-axis root for any $\tau_p\ge 0$. The undelayed Jacobian is Hurwitz ($C_E=-0.850$).
\end{numresult}

\section{Iso-gain sign reversal}

\begin{proposition}[Iso-gain sign flip]
\label{prop:flip}
Replace $C_Z$ by $-C_Z$ in the companion's gated Candidate~A linearisation, leaving every other coefficient unchanged. Then $H$ is identical, the frequencies are identical, and the fundamental delays become
\[
\tau_-=128.374,\qquad \tau_+=70.697
\]
years, which are the companion values shifted by $\pi/\omega$ on each family. Both crossings remain local Hopfs.
\end{proposition}

\begin{proof}
The cubic $H$ is unchanged because it depends on $C_Z$ only through $C_Z^2$. The characteristic identity becomes
\[
P(i\omega)=(-C_Z)L(i\omega)e^{-i\omega\tau}
=C_Z L(i\omega)e^{-i\omega\tau}e^{\pm i\pi}.
\]
If $\tau_0$ solves the identity for $+C_Z$, then $\tau_0+\pi/\omega$ solves it for $-C_Z$. Adding $2\pi k/\omega$ exhausts the lattice. Evaluating the argument formula of~\cite{paperI} on the two positive roots of $H$ at Candidate~A produces the fundamentals $128.374$ and $70.697$.
\end{proof}

The reversed-gain linearisation has loop gain $1.016>1$ and retains the factor $\eta E^*/\Delta_{\mathrm{ref}}$. Definition~\ref{def:prot} changes the modulus as well as the sign.

\section{Two delays}

Let $Z_m$ and $Z_p$ be two copies of the companion filter, or a single $Z$ observed with two delays. The interpolated effort law is
\begin{equation}
\dot E=\Bigl(1-\frac{E}{E_{\max}}\Bigr)
\Bigl[\chi_m F_m\bigl(E,Z(t-\tau_m)\bigr)+\chi_p\eta_p\bigl(E_{\mathrm{cap}}(Z(t-\tau_p))-E\bigr)\Bigr],
\label{eq:two}
\end{equation}
with $\chi_m,\chi_p\ge 0$. The companion is $(\chi_m,\chi_p)=(1,0)$. Pure protection is $(0,1)$.

Linearising~\eqref{eq:two} at an interior rest where both brackets vanish separately (the Candidate~A point, under~\eqref{eq:cap}) gives
\begin{equation}
\dot y=C_E y+C_m z(t-\tau_m)+C_p z(t-\tau_p),
\label{eq:y}
\end{equation}
with $C_m=\chi_m C_Z^{\mathrm{mob}}$, $C_p=\chi_p C_Z^{\mathrm{prot}}$, and $C_E$ the sum of the two gate-adjusted $E$-derivatives.

\begin{theorem}[Two-delay characteristic identity]
\label{thm:char}
The characteristic function of the stock--memory block coupled to~\eqref{eq:y} is
\begin{equation}
P(\lambda)-L(\lambda)\bigl(C_m e^{-\lambda\tau_m}+C_p e^{-\lambda\tau_p}\bigr)=0,
\label{eq:qp}
\end{equation}
where $P(\lambda)=(\lambda-A_N)(\lambda+d)(\lambda-C_E)$ and $L(\lambda)=B_E(\lambda-A_N)+A_E B_N$ are the companion polynomials.
\end{theorem}

\begin{proof}
The variational system is
\[
\dot x=A_N x+A_E y,\qquad
\dot z=B_N x+B_E y-d z,\qquad
\dot y=C_E y+C_m z(t-\tau_m)+C_p z(t-\tau_p).
\]
Its characteristic matrix is
\[
\begin{pmatrix}
\lambda-A_N&0&-A_E\\
-B_N&\lambda+d&-B_E\\
0&-C_m e^{-\lambda\tau_m}-C_p e^{-\lambda\tau_p}&\lambda-C_E
\end{pmatrix}.
\]
Expanding the determinant along the third row produces
\[
(\lambda-C_E)(\lambda-A_N)(\lambda+d)
-\bigl(C_m e^{-\lambda\tau_m}+C_p e^{-\lambda\tau_p}\bigr)\bigl(B_E(\lambda-A_N)+A_E B_N\bigr),
\]
which is $P(\lambda)-L(\lambda)(C_m e^{-\lambda\tau_m}+C_p e^{-\lambda\tau_p})$.
\end{proof}

\begin{theorem}[Weighted small gain]
\label{thm:weighted}
If
\[
\sup_{\omega>0}\frac{\bigl(|C_m|+|C_p|\bigr)\,|L(i\omega)|}{|(i\omega-A_N)(i\omega+d)(i\omega-C_E)|}<1,
\]
then~\eqref{eq:qp} has no imaginary-axis root for any $\tau_m,\tau_p\ge 0$.
\end{theorem}

\begin{proof}
If $\lambda=i\omega$ solves~\eqref{eq:qp}, then
\[
P(i\omega)=L(i\omega)\bigl(C_m e^{-i\omega\tau_m}+C_p e^{-i\omega\tau_p}\bigr),
\]
hence
\[
|P(i\omega)|\le |L(i\omega)|\bigl(|C_m|+|C_p|\bigr),
\]
using $|e^{-i\omega\tau}|=1$. Dividing by $|(i\omega-A_N)(i\omega+d)(i\omega-C_E)|$ yields a loop gain of at least one, contradicting the hypothesis.
\end{proof}

\begin{corollary}[Mobilising weight]
\label{cor:weight}
At Candidate~A, the pure mobilising loop gain exceeds $1$ and the pure protective loop gain is $0.080$. Assume the common equilibrium and all linear coefficients depend continuously on $\chi_m$, the characteristic denominator remains nonzero on the imaginary axis, and the protective endpoint has a strict gain margin. Then there exists $\chi_m^*\in(0,1)$ such that every interpolation with $\chi_m<\chi_m^*$ and $\chi_p=1-\chi_m$ satisfies Theorem~\ref{thm:weighted}. A Hopf of the interpolated system therefore requires a sufficiently large mobilising weight; it cannot be produced by decreasing $\tau_p$ alone.
\end{corollary}

\section{Sampled protection}

The companion's sample-and-hold monodromy $M(T_r)$ is the review operator for a continuously linearised effort law. Replacing $C_Z$ by $C_Z^{\mathrm{prot}}<0$ in that formula yields a different matrix $M_p(T_r)$.

\begin{proposition}[Annual review under protection]
\label{prop:annual}
With the gains~\eqref{eq:CZ}, $\rho\bigl(M_p(1)\bigr)=0.9838<1$. Annual review of the quota-tracking channel is linearly stable at Candidate~A. The same Euler hold map crosses $\rho=1$ at $T_r=2.306$; that crossing is a discretisation of the explicit Euler step $1+T_r C_E$ with $C_E=-0.850$, not a Hopf of the protective DDE (Theorem~\ref{thm:nohopf}).
\end{proposition}

\begin{proof}
Between reviews, $E$ is frozen, so the variational $(x,z,y)$-system is $\dot\xi=A_{\mathrm{hold}}\xi$ with
\[
A_{\mathrm{hold}}=\begin{pmatrix}A_N&0&A_E\\ B_N&-d&B_E\\ 0&0&0\end{pmatrix}.
\]
The review update linearises to the shear with $C_E,C_Z$ from~\eqref{eq:CZ}. The monodromy is therefore
\[
M_p(T_r)=\begin{pmatrix}1&0&0\\0&1&0\\0&T_r C_Z&1+T_r C_E\end{pmatrix}\exp(A_{\mathrm{hold}}T_r),
\]
the same algebraic expression as in~\cite{paperI} with the protective gains. Evaluating the eigenvalues of this $3\times 3$ matrix at $T_r=1$ gives spectral radius $0.9838$. On the grid $T_r\in[0.2,20]$ the spectral radius is strictly less than one on $[0.2,2.306)$ and strictly greater than one on $(2.306,20]$, so a unique crossing occurs at $T_r=2.306$. That crossing is produced by the Euler factor $1+T_r C_E$ with $C_E=-0.850$, not by an imaginary root of the continuous characteristic function (Theorem~\ref{thm:nohopf}).
\end{proof}

At Candidate~A the mobilising hold map of~\cite{paperI} is unstable for $T_r=1$, because the undelayed mobilising Jacobian is already unstable. The protective map $M_p$ does not inherit that instability.

\section{Pacing}

\begin{theorem}[Channel-specific pacing]
\label{thm:policy}
For~\eqref{eq:core} with the mobilising bracket of~\cite{paperI}, the equilibrium is linearly unstable for $0<\tau<\tau_-$. For~\eqref{eq:prot} at Candidate~A, the equilibrium is linearly stable for every $\tau_p\ge 0$. For~\eqref{eq:two} any delay-induced instability, if it occurs, lies in a region of $(\tau_m,\chi_m)$ and is independent of $\tau_p$ whenever Theorem~\ref{thm:weighted} applies.
\end{theorem}

\begin{proof}
Theorem~\ref{thm:nohopf} and Corollary~\ref{cor:weight}. A root of~\eqref{eq:qp} at the mobilising $\tau_-$ requires $C_m$ large enough that the weighted gain reaches one.
\end{proof}

\section{Discussion}

Equation~\eqref{eq:prot} has no local Hopf at Candidate~A, so it has no small periodic branch born from that equilibrium. Roots of~\eqref{eq:qp} for $\chi_m>\chi_m^*$ may be continued in $(\tau_m,\tau_p)$.

\section{Conclusion}

Linear delay-induced instability of the three-state core requires a positive delayed gain of sufficient modulus. Quota tracking at the Candidate~A stock equilibrium has negative gain, loop gain $0.080$, and no Hopf. The two-channel characteristic identity separates the two mechanisms.

\begin{thebibliography}{9}
\bibitem{paperI}
Y.\ Name, ``Scarcity-driven capital liquidation and delay-amplified instability: a vector-valued flow-balance framework for sustainability,'' companion manuscript, 2026.
\end{thebibliography}

\end{document}
